\section{微扰 QCD 与树级过程}\label{sec:qcd-tree-processes-dp10}

上一节已经给出 QCD 作用量。对刚接触图规则的读者，最重要的不是先背图，而是知道图中的每个对象怎样由作用量产生。把作用量在真空附近按场的次数展开：二次项定义自由线性方程，其 Fourier 核的逆就是\emph{传播子}；三次、四次项分别含三个、四个场因子，对这些场作泛函微分就得到相应\emph{顶角因子}。再把顶角与外部在壳旋量、偏振矢量相连，并对内部线携带的未定四动量积分，才得到具体散射振幅。

一个 Feynman 图因此只是 Dyson--Wick 展开中某一组收缩的记号：外线指定渐近态，内部线记录传播子的动量核，顶点记录相互作用项。\emph{树级}（tree level）指没有独立圈动量积分的最低非零微扰阶；\emph{一圈}（one loop）则至少含一个独立 $\int\dd^4k$。颜色生成元的非交换性会同时进入树图颜色因子和圈图群论系数。不同味夸克散射最适合作为第一个例子，因为它只有一条最低阶交换拓扑，可以先把 Dirac 自旋代数与颜色代数分开计算；湮灭过程和含外部胶子的过程随后只是在这两类代数之上增加新的拓扑与顶角。

\subsection{QCD 图规则}

外态归一化、整体四动量守恒 $\delta$ 函数和 SI 相空间前因子采用统一散射约定；图规则采用无量纲约化不变量振幅 $\mathcal M$，内部线与顶角如下。取 Feynman 规范 $\xi=1$。夸克和胶子的物理四动量分别记为 $p^\mu$ 与 $k^\mu$，均具有 $\mathrm{kg\,m\,s^{-1}}$ 量纲。

夸克二次核的逆给出
\begin{equation}
 \boxed{
 S_f(p)
 =\frac{\ii(\slashed p+m_fc)}{p^2-m_f^2c^2+\ii0^+}\,\delta_{ij},}
 \label{eq:qcd-quark-propagator-dp10}
\end{equation}
而 Feynman 规范下胶子二次核的逆为
\begin{equation}
 \boxed{
 D_{\mu\nu}^{ab}(k)
 =\frac{-\ii\delta^{ab}\eta_{\mu\nu}}
 {k^2+\ii0^+}.}
 \label{eq:qcd-gluon-propagator-dp10}
\end{equation}
这里 $i,j$ 是基本表示颜色指标，$a,b$ 是伴随表示指标。

夸克--胶子相互作用给出约化顶角
\begin{equation}
 \boxed{
 V_{q\bar q g}^{a\mu}=-\ii g_s\gamma^\mu T^a.}
 \label{eq:qcd-quark-gluon-vertex-dp10}
\end{equation}
这个顶角同时作用在两个彼此独立的空间上：$\gamma^\mu$ 作用在 Dirac 旋量指标，$T^a$ 作用在夸克颜色指标。若入射夸克颜色为 $j$、出射夸克颜色为 $i$，则颜色振幅就是 $(T^a)_{ij}$；伴随指标 $a$ 指明与该顶角相连的胶子颜色方向。所谓“胶子把夸克颜色从 $j$ 变到 $i$”正是这个矩阵元的物理含义，而不是额外附加在图上的规则。

整体相位由全书统一的
$S=\Tord\exp[(\ii/\hbar c)\int\dd^4 x\,\Lag_{\rm int}]$
约定固定；所有过程必须在同一约定下相干求和。

\begin{derivation}[从\eqnrefs{eq:qcd-lagrangian-si-dp10}到\eqnrefs{eq:qcd-quark-gluon-vertex-dp10}]
把协变导数
\[
D_\mu=\partial_\mu+\ii g_s\mathscr C_\mu^aT^a
\]
代入夸克动能项，先把自由项与相互作用项严格分开：
\begin{align*}
\bar\psi_f\ii\hbar c\gamma^\mu D_\mu\psi_f
&=\bar\psi_f\ii\hbar c\gamma^\mu\partial_\mu\psi_f
+\bar\psi_f\ii\hbar c\gamma^\mu
(\ii g_s\mathscr C_\mu^aT^a)\psi_f\\
&=\bar\psi_f\ii\hbar c\gamma^\mu\partial_\mu\psi_f
-\hbar c g_s\bar\psi_f\gamma^\mu T^a\psi_f\,\mathscr C_\mu^a.
\end{align*}
因此三场相互作用拉氏量就是
\[
\Lag_{q\bar q g}
=-\hbar c g_s\bar\psi_{f,i}\gamma^\mu(T^a)_{ij}\psi_{f,j}\mathscr C_\mu^a.
\]
Dyson 展开的一阶项为
\begin{align*}
S_{q\bar q g}^{(1)}
&=\frac{\ii}{\hbar c}\int\dd^4x\,\Lag_{q\bar q g}(x)\\
&=-\ii g_s\int\dd^4x\,
\bar\psi_{f,i}(x)\gamma^\mu(T^a)_{ij}\psi_{f,j}(x)\mathscr C_\mu^a(x).
\end{align*}
把三条场都作 Fourier 展开，局域乘积中的时空积分产生整体四动量守恒因子
\[
\int\dd^4x\,
\exp\!\left[-\frac{\ii}{\hbar}
(p_1+p_2+k)\cdot x\right]
=(2\pi\hbar)^4\delta^{(4)}(p_1+p_2+k),
\]
这里按图规则约定三条四动量全部流入顶角。剥去所有图共有的动量守恒 $\delta$ 函数以后，剩下的局域矩阵系数恰为
\[
\boxed{-\ii g_s\gamma^\mu(T^a)_{ij}}.
\]
因此一个 QCD 顶角的三部分来源可以逐项追溯：$-\ii g_s$ 来自 Dyson 指数与协变导数中的耦合，$\gamma^\mu$ 来自 Dirac 电流，$T^a$ 来自颜色联络在基本表示上的生成元。图规则只是把这三个作用量系数压缩成一个可重复使用的记号。
\end{derivation}

三条胶子动量 $p,q,r$ 全部取为流入顶角并满足 $p+q+r=0$ 时，三胶子顶角为
\begin{equation}
 \boxed{
 V_{3g}^{abc,\mu\nu\rho}(p,q,r)
 =-g_sf^{abc}\Big[
 \eta^{\mu\nu}(p-q)^\rho
 +\eta^{\nu\rho}(q-r)^\mu
 +\eta^{\rho\mu}(r-p)^\nu
 \Big].}
 \label{eq:qcd-three-gluon-vertex-dp10}
\end{equation}
动量差来自 Yang--Mills 三次项中的导数，而 $f^{abc}$ 保留了非阿贝尔颜色结构。四胶子顶角不含额外动量因子；对全部流入的指标 $(a,\mu),(b,\nu),(c,\rho),(d,\sigma)$，
\begin{align*}
 V_{4g}^{abcd,\mu\nu\rho\sigma}
 =-\ii g_s^2\Big[
 &f^{abe}f^{cde}(\eta^{\mu\rho}\eta^{\nu\sigma}-\eta^{\mu\sigma}\eta^{\nu\rho})\\
 +&f^{ace}f^{bde}(\eta^{\mu\nu}\eta^{\rho\sigma}-\eta^{\mu\sigma}\eta^{\nu\rho})\\
 +&f^{ade}f^{bce}(\eta^{\mu\nu}\eta^{\rho\sigma}-\eta^{\mu\rho}\eta^{\nu\sigma})
 \Big].
\end{align*}
改变“全部流入”的动量方向或结构常数的指标顺序会改变某些中间符号；可观测量要求的是同一作用量约定下全部图的一致组合。

规范固定还产生鬼场传播子
\begin{equation*}
 D_{\rm gh}^{ab}(k)=\frac{\ii\delta^{ab}}{k^2+\ii0^+},
\end{equation*}
以及由 $-\hbar c g_sf^{abc}(\partial^\mu\bar c^a)\mathscr C_\mu^b c^c$ 直接读出的鬼--反鬼--胶子顶角。若反鬼四动量 $k^\mu$ 流入顶角，约化顶角因子为 $-g_sf^{abc}k^\mu$（整体 $\ii$ 相位依全书统一 $S$ 矩阵约定固定）。鬼场只出现在内部线中，不属于物理渐近态；它在圈图中补偿协变量子化引入的非物理规范自由度。

\begin{derivation}[从\eqnrefs{eq:qcd-gluon-propagator-dp10}到\eqnrefs{eq:qcd-three-gluon-vertex-dp10}]
Feynman 规范 $\xi=1$ 下，把纯规范二次项与规范固定项分部积分后写成
\begin{equation*}
 \mathcal L_2
 =\frac{\hbar c}{2}\,
 \mathscr C_\mu^a\,\eta^{\mu\nu}\partial^2\mathscr C_\nu^a .
\end{equation*}
因此作用在规范场上的二次微分算符是
\begin{equation*}
 (\mathsf K\mathscr C)_\mu^a
 =\delta^{ab}\eta_{\mu\nu}\partial^2\mathscr C^{b\nu}.
\end{equation*}
采用
\begin{equation*}
 \mathscr C_\mu^a(x)
 =\int\frac{\dd^4p}{(2\pi\hbar)^4}
 \widetilde{\mathscr C}_\mu^a(p)
 \ee^{-\ii p\cdot x/\hbar},
\end{equation*}
有 $\partial^2\mapsto-p^2/\hbar^2$，故
\begin{equation*}
 \mathsf K_{\mu\nu}^{ab}(p)
 =-\frac{p^2}{\hbar^2}\delta^{ab}\eta_{\mu\nu}.
\end{equation*}
Feynman 边界条件把 $p^2$ 替换为 $p^2+\ii0^+$。定义逆核 $\mathsf K_F^{-1}$ 使
\begin{equation*}
 \mathsf K_{\mu\rho}^{ac}(p)
 (\mathsf K_F^{-1})^{cb\rho}{}_{\nu}(p)
 =\delta^{ab}\delta_\mu{}^\nu,
\end{equation*}
于是逐指标求逆得到
\begin{equation*}
 (\mathsf K_F^{-1})_{\mu\nu}^{ab}(p)
 =-\hbar^2\frac{\delta^{ab}\eta_{\mu\nu}}{p^2+\ii0^+}.
\end{equation*}
在统一的约化振幅约定中，玻色 Feynman 线等于 $(\ii/\hbar^2)\mathsf K_F^{-1}$，从而
\begin{equation*}
 D_{\mu\nu}^{ab}(p)
 =\frac{\ii}{\hbar^2}(\mathsf K_F^{-1})_{\mu\nu}^{ab}(p)
 =\frac{-\ii\delta^{ab}\eta_{\mu\nu}}{p^2+\ii0^+},
\end{equation*}
即正文\eqnrefs{eq:qcd-gluon-propagator-dp10}。

三胶子规则来自 Yang--Mills 作用量的三次部分
\begin{equation*}
 \mathcal L_{3g}
 =\hbar c g_s f^{abc}
 (\partial_\mu\mathscr C_\nu^a)
 \mathscr C^{b\mu}\mathscr C^{c\nu}.
\end{equation*}
令三条流入场分别为 $(a,\alpha,p)$、$(b,\beta,q)$、$(c,\gamma,r)$，并约定
\begin{equation*}
 p+q+r=0.
\end{equation*}
三次场都是同一个玻色规范场，因此对作用量取三次泛函导数时，导数场可以对应三条外腿中的任意一条。把所有排列写到统一的颜色次序 $f^{abc}$ 后，三个 Lorentz 系数分别为
\begin{align*}
 \mathscr V_1^{\alpha\beta\gamma}
 &=\eta^{\alpha\beta}(p-q)^\gamma,\\
 \mathscr V_2^{\alpha\beta\gamma}
 &=\eta^{\beta\gamma}(q-r)^\alpha,\\
 \mathscr V_3^{\alpha\beta\gamma}
 &=\eta^{\gamma\alpha}(r-p)^\beta.
\end{align*}
这里每一个动量差都来自两种外腿排列的差：交换相邻两个胶子场时，$f^{abc}$ 变号，而 Fourier 导数把相应流入动量带入同一个 Lorentz 缩并。于是三阶作用量系数给出完整约化顶角
\begin{equation*}
 \boxed{
 V_{3g}^{abc,\alpha\beta\gamma}(p,q,r)
 =-g_sf^{abc}
 \left[
 \eta^{\alpha\beta}(p-q)^\gamma
 +\eta^{\beta\gamma}(q-r)^\alpha
 +\eta^{\gamma\alpha}(r-p)^\beta
 \right],}
\end{equation*}
即正文\eqnrefs{eq:qcd-three-gluon-vertex-dp10}。交换任意两条完整外腿，例如 $(a,\alpha,p)\leftrightarrow(b,\beta,q)$，$f^{abc}$ 与方括号同时变号，因此完整顶角保持玻色交换对称。
\end{derivation}


\subsection{不同味的 $qq'\to qq'$}

QCD 树图里同时出现自旋迹与颜色代数。最适合第一次把两者分开的过程，是两个不同味夸克的弹性散射：味不同排除了交换同种费米子的 $u$ 道图，只剩单胶子 $t$ 道。先保留两个夸克质量
\begin{equation*}
 q(p_1,i;m_1)+q'(p_2,j;m_2)
 \longrightarrow
 q(p_3,k;m_1)+q'(p_4,l;m_2),
\end{equation*}
并定义
\begin{equation*}
 s=(p_1+p_2)^2,
 \qquad
 t=(p_1-p_3)^2,
 \qquad
 u=(p_1-p_4)^2.
\end{equation*}
外腿质量壳条件给出严格恒等式
\begin{equation}
 \boxed{s+t+u=2(m_1^2+m_2^2)c^2.}
 \label{eq:qqprime-mandelstam-massive-dp115}
\end{equation}
因此无质量关系 $s+t+u=0$ 只在最后取 $m_1,m_2\to0$ 时出现。

\begin{figure}[htbp]
 \centering
 \begin{tikzpicture}[
 x=1cm,y=1cm,
 gluon/.style={draw,line width=0.75pt,decorate,decoration={coil,aspect=0.65,amplitude=1.0pt,segment length=4.2pt}}
]
% qq' scattering
\begin{scope}[shift={(0,0)}]
 \node[font=\bfseries] at (0,1.7) {(a) $qq'\to qq'$};
 \coordinate (v1) at (0,0.55); \coordinate (v2) at (0,-0.55);
 \draw[fermion] (-1.8,1.0)--(v1)--(1.8,1.0);
 \draw[fermion] (-1.8,-1.0)--(v2)--(1.8,-1.0);
 \draw[gluon] (v1)--(v2);
 \fill[fvertex] (v1); \fill[fvertex] (v2);
 \node[font=\small] at (0.45,0) {$g^a$};
\end{scope}
% qqbar -> gg, t u s
\begin{scope}[shift={(5.6,0)}]
 \node[font=\bfseries] at (0,1.7) {(b) $q\bar q\to gg$: $t+u+s$};
 \begin{scope}[shift={(-1.6,0)},scale=.67]
  \coordinate (v1) at (-0.45,0); \coordinate (v2) at (0.45,0);
  \draw[fermion] (-1.35,0.85)--(v1)--(v2);
  \draw[antifermion] (-1.35,-0.85)--(v2);
  \draw[gluon] (v1)--(1.35,0.85);
  \draw[gluon] (v2)--(1.35,-0.85);
  \fill[fvertex] (v1);\fill[fvertex] (v2);
  \node at (0,-1.25) {$t$};
 \end{scope}
 \node[font=\large] at (0,0) {$+$};
 \begin{scope}[shift={(1.6,0)},scale=.67]
  \coordinate (v1) at (-0.45,0); \coordinate (v2) at (0.45,0);
  \draw[fermion] (-1.35,0.85)--(v1)--(v2);
  \draw[antifermion] (-1.35,-0.85)--(v2);
  \draw[gluon] (v2)--(1.35,0.85);
  \draw[gluon] (v1)--(1.35,-0.85);
  \fill[fvertex] (v1);\fill[fvertex] (v2);
  \node at (0,-1.25) {$u$};
 \end{scope}
 \node[font=\large] at (3.0,0) {$+$};
 \begin{scope}[shift={(4.4,0)},scale=.67]
  \coordinate (v1) at (-0.45,0); \coordinate (v2) at (0.45,0);
  \draw[fermion] (-1.35,0.85)--(v1);
  \draw[antifermion] (-1.35,-0.85)--(v1);
  \draw[gluon] (v1)--(v2);
  \draw[gluon] (v2)--(1.35,0.85);
  \draw[gluon] (v2)--(1.35,-0.85);
  \fill[fvertex] (v1);\fill[fvertex] (v2);
  \node at (0,-1.25) {$s$};
 \end{scope}
\end{scope}
\end{tikzpicture}

 \caption{左：不同味的 $qq'\to qq'$ 只有单胶子 $t$ 道交换。右：$q\bar q\to gg$ 必须把夸克交换的两种排序与非阿贝尔 $s$ 道三胶子图相干相加，Ward 恒等式才成立。}
 \label{fig:qcd-tree-processes-dp10}
\end{figure}

\figref{fig:qcd-tree-processes-dp10} 左图对应本小节唯一的 $t$ 道收缩；右图将在后面的 $q\bar q\to gg$ 中使用，其中两种夸克线排序和三胶子图必须在振幅层相干相加。

沿两条费米子线写振幅，并省去所有图共有的整体相位，得到
\begin{equation}
 \boxed{
 \mathcal M_{qq'}
 =\frac{g_s^2}{t}
 \bigl[\bar u(p_3)\gamma^\mu T^a_{ki}u(p_1)\bigr]
 \bigl[\bar u(p_4)\gamma_\mu T^a_{lj}u(p_2)\bigr].}
 \label{eq:qqprime-amplitude-dp10}
\end{equation}
颜色指标和自旋指标作用在不同张量因子上，因此平方振幅可以把两类求和分别完成。记
\begin{equation*}
 \Sigma_{12}\equiv(m_1^2+m_2^2)c^2 .
\end{equation*}
对一般 $SU(N_c)$ 且 $T_F=1/2$，初态颜色平均给 $(N_c^2-1)/(4N_c^2)$；完整有限质量 Dirac 迹给出
\begin{equation}
 \boxed{
 \overline{|\mathcal M(qq'\to qq')|^2}
 =g_s^4\frac{N_c^2-1}{2N_c^2}
 \frac{(s-\Sigma_{12})^2+(u-\Sigma_{12})^2+2t\Sigma_{12}}{t^2}.}
 \label{eq:qqprime-M2-massive-dp115}
\end{equation}
连续取 $m_1,m_2\to0$ 后才得到
\begin{equation}
 \boxed{
 \overline{|\mathcal M(qq'\to qq')|^2}_{m_1=m_2=0}
 =g_s^4\frac{N_c^2-1}{2N_c^2}
 \frac{s^2+u^2}{t^2}.}
 \label{eq:qqprime-M2-final-dp10}
\end{equation}
物理 QCD 再取 $N_c=3$，颜色系数为 $4/9$。

以下角分布专门讨论无质量子模型。在质心系中
\begin{equation*}
 t=-\frac s2(1-\cos\theta),
 \qquad
 u=-\frac s2(1+\cos\theta).
\end{equation*}
代入统一的二体相空间公式并用 $g_s^2=4\pi\alpha_s$，得到
\begin{equation}
 \boxed{
 \frac{\dd\sigma(qq'\to qq')}{\dd\Omega}
 =\frac{\hbar^2\alpha_s^2}{8s}
 \frac{N_c^2-1}{N_c^2}
 \frac{4+(1+\cos\theta)^2}{(1-\cos\theta)^2}.}
 \label{eq:qqprime-angular-cross-section-dp10}
\end{equation}
对 $N_c=3$，前因子化为 $\hbar^2\alpha_s^2/(9s)$。\eqnrefs{eq:qqprime-angular-cross-section-dp10} 可分成运动学形状与整体尺度。定义
\begin{equation}
 K_{qq'}(\theta)
 \equiv
 \frac{4+(1+\cos\theta)^2}{(1-\cos\theta)^2}.
 \label{eq:qcd-qqprime-angular-kernel-dp68}
\end{equation}
\figref{fig:qcd-qqprime-angular-dp68} 画出 $K_{qq'}(\theta)/K_{qq'}(90^\circ)$。对数纵轴把前向增强直接显示为随 $\theta\to0$ 的快速增长；曲线逐点由正文树级截面的角因子计算，没有引入拟合参数。

\begin{figure}[htbp]
 \centering
 \includegraphics[width=0.72\textwidth]{qcd_qqprime_angular_dp68.pdf}
 \caption{由\eqnrefs{eq:qcd-qqprime-angular-kernel-dp68} 直接计算的不同味 $qq'\to qq'$ 树级角分布核，并以 $90^\circ$ 的值归一化。前向增强来自无质量 $t$ 道胶子传播子的 $1/t^2$ 结构。}
 \label{fig:qcd-qqprime-angular-dp68}
\end{figure}

$\theta\to0$ 的前向增强来自无质量 $t$ 道传播子的长程奇异性，与 QED 的 Coulomb 交换具有相同的运动学来源。图中的夸克、胶子外腿与实验直接看到的强子属于不同描述层级。

在足够大的动量转移下，碰撞的短时间、短距离部分可用近似自由的夸克或胶子线描述；参与这种硬子过程的夸克、反夸克和胶子统称为\emph{部分子}（parton）。长距离强相互作用又把有色部分子组织成无色强子，所以探测器不能孤立读取单个部分子。

高能末态部分子先发生软--共线 QCD 辐射，随后在强耦合长距离区转化成颜色单态强子；这个从有色部分子描述过渡到实际强子末态的非微扰过程称为\emph{强子化}（hadronization）。实验中不会把某个理论部分子与某个单独强子一一对应，而是用指定的喷注算法把方向接近的一组末态强子聚类成\emph{喷注}（jet）。因此“一个 jet”首先是一个红外安全的测量定义，其能量和方向在高能极限下追踪产生它的硬部分子。有限的强子喷注截面必须把短程部分子振幅、强子长程结构、软--共线辐射以及喷注算法定义放在同一因子化计算中。

\begin{derivation}[从\eqnrefs{eq:qqprime-amplitude-dp10}到\eqnrefs{eq:qqprime-M2-massive-dp115}]
从\eqnrefs{eq:qqprime-amplitude-dp10} 出发，初态自旋平均和颜色平均可分开写成
\begin{equation*}
 \overline{|\mathcal M|^2}
 =\frac{g_s^4}{t^2}\,
 \mathcal S_{\mu\nu}^{(1)}\mathcal S_{(2)}^{\mu\nu}\,\mathcal C,
\end{equation*}
其中
\begin{align*}
 \mathcal S_{\mu\nu}^{(1)}
 &=\frac12\Tr[(\slashed p_3+m_1c)\gamma_\mu(\slashed p_1+m_1c)\gamma_\nu],\\
 \mathcal S_{(2)}^{\mu\nu}
 &=\frac12\Tr[(\slashed p_4+m_2c)\gamma^\mu(\slashed p_2+m_2c)\gamma^\nu].
\end{align*}
颜色部分与质量无关：
\begin{align*}
 \mathcal C
 &=\frac1{N_c^2}\sum_{i,j,k,l}\sum_{a,b}
 (T^a)_{ki}(T^a)_{lj}(T^b)_{ik}(T^b)_{jl}\\
 &=\frac1{N_c^2}\sum_{a,b}\Tr(T^aT^b)\Tr(T^aT^b)\\
 &=\frac{T_F^2(N_c^2-1)}{N_c^2}
 =\frac{N_c^2-1}{4N_c^2}.
\end{align*}
Dirac 迹恒等式给
\begin{align*}
 \mathcal S_{\mu\nu}^{(1)}
 &=2\left[p_{3\mu}p_{1\nu}+p_{3\nu}p_{1\mu}
 -\eta_{\mu\nu}(p_1\!\cdot p_3-m_1^2c^2)\right],\\
 \mathcal S_{(2)}^{\mu\nu}
 &=2\left[p_4^{\mu}p_2^{\nu}+p_4^{\nu}p_2^{\mu}
 -\eta^{\mu\nu}(p_2\!\cdot p_4-m_2^2c^2)\right].
\end{align*}
先在四矢量标量积层面完成张量收缩。把两个方括号分别记为 $A_{\mu\nu}$ 与 $B^{\mu\nu}$，则逐项展开 $A_{\mu\nu}B^{\mu\nu}$ 后，所有
$(p_1\!\cdot p_3)(p_2\!\cdot p_4)$ 的纯迹项相消，得到
\begin{align*}
\mathcal S_{\mu\nu}^{(1)}\mathcal S_{(2)}^{\mu\nu}
=8\Big[&
(p_3\!\cdot p_4)(p_1\!\cdot p_2)
+(p_3\!\cdot p_2)(p_1\!\cdot p_4)\\
&-m_2^2c^2(p_1\!\cdot p_3)
-m_1^2c^2(p_2\!\cdot p_4)
+2m_1^2m_2^2c^4\Big].
\end{align*}
现在才把这些内积换成 Mandelstam 变量。除
\begin{align*}
 2p_1\!\cdot p_2&=s-\Sigma_{12},
&2p_1\!\cdot p_4&=\Sigma_{12}-u,\\
 2p_1\!\cdot p_3&=2m_1^2c^2-t,
&2p_2\!\cdot p_4&=2m_2^2c^2-t,
\end{align*}
外，由四动量守恒还可写成
\begin{equation*}
2p_3\!\cdot p_4=s-\Sigma_{12},
\qquad
2p_3\!\cdot p_2=\Sigma_{12}-u.
\end{equation*}
前两项因而给
\begin{equation*}
\frac14\left[(s-\Sigma_{12})^2+(u-\Sigma_{12})^2\right],
\end{equation*}
而三个显式质量项组合为
\begin{align*}
&-m_2^2c^2\left(m_1^2c^2-\frac t2\right)
-m_1^2c^2\left(m_2^2c^2-\frac t2\right)
+2m_1^2m_2^2c^4\\
&\qquad=\frac{t\Sigma_{12}}2.
\end{align*}
乘回最外层因子 $8$，得到
\begin{equation*}
 \mathcal S_{\mu\nu}^{(1)}\mathcal S_{(2)}^{\mu\nu}
 =2\left[(s-\Sigma_{12})^2+(u-\Sigma_{12})^2+2t\Sigma_{12}\right].
\end{equation*}
与颜色因子相乘即得\eqnrefs{eq:qqprime-M2-massive-dp115}。令 $m_1,m_2\to0$ 后，$\Sigma_{12}\to0$ 且 $s+t+u\to0$，从而得到\eqnrefs{eq:qqprime-M2-final-dp10}。
\end{derivation}

\subsection{湮灭过程与颜色多重度}

交叉关系把同一个夸克--胶子顶角组织成不同的散射通道。考虑不同味的
$q\bar q\to Q\bar Q$，先保留初态质量 $m_q$ 与末态质量 $m_Q$；树级只有 $s$ 道胶子交换。由两端夸克--胶子顶角和中间胶子传播子组成的树级振幅为
\begin{equation}
\boxed{
 \mathcal M(q\bar q\to Q\bar Q)
 =\frac{g_s^2}{s}
 [\bar v(p_2)\gamma^\mu T^a u(p_1)]
 [\bar u(p_3)\gamma_\mu T^a v(p_4)].}
\label{eq:qqbar-to-QQbar-amplitude-dp68}
\end{equation}
对初态自旋和颜色平均、对末态求和后，记
\begin{equation*}
 \Sigma_{qQ}\equiv(m_q^2+m_Q^2)c^2,
\end{equation*}
一般 $SU(N_c)$ 的有限质量结果为
\begin{equation}
 \boxed{
 \overline{|\mathcal M(q\bar q\to Q\bar Q)|^2}
 =g_s^4\frac{N_c^2-1}{2N_c^2}
 \frac{
 (t-\Sigma_{qQ})^2+(u-\Sigma_{qQ})^2
 +2s\Sigma_{qQ}
 }{s^2}.}
 \label{eq:qqbar-to-QQbar-M2-dp10}
\end{equation}
外腿质量壳条件给
\[
s+t+u=2(m_q^2+m_Q^2)c^2=2\Sigma_{qQ}.
\]
只有在 $m_q,m_Q\to0$ 的连续极限中才退化为
\[
\overline{|\mathcal M|^2}
\longrightarrow
g_s^4\frac{N_c^2-1}{2N_c^2}
\frac{t^2+u^2}{s^2}.
\]
这样同一个公式同时覆盖轻味产生与重味阈值以上的树级部分子过程；计算有限质量截面时，$m_Q$ 还会通过二体相空间速度
\(\beta_Q=\sqrt{1-4m_Q^2c^2/s}\)
进入末态相空间。

颜色多重度在电磁产生中更直接。电磁流在颜色指标上严格对角，单个光子--夸克顶角的颜色张量是 $\delta_{ij}$。因此末态颜色求和给出
\begin{equation*}
 \sum_{i,j=1}^{N_c}\delta_{ij}\delta_{ij}=N_c.
\end{equation*}
在无质量树级子模型中，$e^+e^-\to\gamma^*\to q_f\bar q_f$ 的 Dirac 迹与 $e^+e^-\to\mu^+\mu^-$ 相同，只多出电荷因子 $Q_f^2$ 与上述颜色迹。于是单个开放夸克味满足
\begin{equation}
\boxed{
 \sigma(e^+e^-\to q_f\bar q_f)
 =N_cQ_f^2\,\sigma(e^+e^-\to\mu^+\mu^-).}
\label{eq:R-ratio-single-flavor-dp68}
\end{equation}
定义包容强子产生比
\begin{equation*}
 R(s)
 =\frac{\sigma(e^+e^-\to\text{强子})}
 {\sigma(e^+e^-\to\mu^+\mu^-)},
\end{equation*}
在把活跃夸克味定义为该能量下允许产生、并对这些活跃味采用无质量树级子模型时，得到
\begin{equation}
 \boxed{
 R_0(s)=N_c\sum_{f\,\text{开放}}Q_f^2.}
 \label{eq:R-ratio-tree-dp10}
\end{equation}
这条关系把不可直接观测的颜色简并度转化为包容截面的台阶结构：当质心能量跨过新的夸克产生阈值，可参与求和的味数增加。强相互作用的次领阶修正不会改变 $N_cQ_f^2$ 的来源，而是在每个近似无质量的活跃味上加入红外安全的有限校正。

\begin{derivation}[从\eqnrefs{eq:qqbar-to-QQbar-amplitude-dp68}到\eqnrefs{eq:qqbar-to-QQbar-M2-dp10}]
同时保留初态质量 $m_q$ 与末态质量 $m_Q$。颜色与 Dirac 空间彼此独立，因此
\begin{equation*}
 \overline{|\mathcal M|^2}
 =\frac{g_s^4}{s^2}\,\mathcal C\,\mathcal S.
\end{equation*}
颜色平均为
\begin{align*}
\mathcal C
&=\frac1{N_c^2}
\sum_{i,j,k,l}\sum_{a,b}
(T^a)_{ji}(T^a)_{kl}
(T^b)_{ij}(T^b)_{lk}\\
&=\frac1{N_c^2}
\sum_{a,b}\Tr(T^aT^b)\Tr(T^aT^b)\\
&=\frac{T_F^2(N_c^2-1)}{N_c^2}
=\frac{N_c^2-1}{4N_c^2},
\end{align*}
其中基本表示取 $T_F=1/2$。

自旋平均后的两个 Dirac 迹为
\begin{align*}
 L^{\mu\nu}
 &\equiv\Tr[(\slashed p_2-m_qc)\gamma^\mu
 (\slashed p_1+m_qc)\gamma^\nu]\\
 &=4\left[p_2^\mu p_1^\nu+p_2^\nu p_1^\mu
 -\eta^{\mu\nu}(p_1\!\cdot p_2+m_q^2c^2)\right],\\
 H_{\mu\nu}
 &\equiv\Tr[(\slashed p_3+m_Qc)\gamma_\mu
 (\slashed p_4-m_Qc)\gamma_\nu]\\
 &=4\left[p_{3\mu}p_{4\nu}+p_{3\nu}p_{4\mu}
 -\eta_{\mu\nu}(p_3\!\cdot p_4+m_Q^2c^2)\right].
\end{align*}
于是
\begin{equation*}
 \mathcal S=\frac14L^{\mu\nu}H_{\mu\nu}.
\end{equation*}
先不换成 $s,t,u$，直接把两个对称张量逐项收缩。记
$a\equiv p_1\!\cdot p_2$、$b\equiv p_3\!\cdot p_4$，得到
\begin{align*}
\mathcal S
=8\Big[&
(p_1\!\cdot p_3)(p_2\!\cdot p_4)
+(p_1\!\cdot p_4)(p_2\!\cdot p_3)\\
&+m_Q^2c^2 a+m_q^2c^2 b
+2m_q^2m_Q^2c^4\Big].
\end{align*}
质量壳条件给
\begin{equation*}
a=\frac{s-2m_q^2c^2}{2},
\qquad
b=\frac{s-2m_Q^2c^2}{2}.
\end{equation*}
所以后三项不是独立的三种质量修正，而是精确合并成
\begin{align*}
&m_Q^2c^2\frac{s-2m_q^2c^2}{2}
+m_q^2c^2\frac{s-2m_Q^2c^2}{2}
+2m_q^2m_Q^2c^4\\
&\qquad=\frac{s}{2}(m_q^2+m_Q^2)c^2.
\end{align*}
因此
\begin{equation*}
 \mathcal S
 =8\left[
 (p_1\!\cdot p_3)(p_2\!\cdot p_4)
 +(p_1\!\cdot p_4)(p_2\!\cdot p_3)
 +\frac{s}{2}(m_q^2+m_Q^2)c^2
 \right].
\end{equation*}
令
\begin{equation*}
 \Sigma_{qQ}\equiv(m_q^2+m_Q^2)c^2.
\end{equation*}
由有限质量 Mandelstam 恒等式
\begin{align*}
 2p_1\!\cdot p_3&=\Sigma_{qQ}-t,
 &2p_2\!\cdot p_4&=\Sigma_{qQ}-t,\\
 2p_1\!\cdot p_4&=\Sigma_{qQ}-u,
 &2p_2\!\cdot p_3&=\Sigma_{qQ}-u,
\end{align*}
可得
\begin{equation*}
 \mathcal S
 =2\left[(t-\Sigma_{qQ})^2
 +(u-\Sigma_{qQ})^2
 +2s\Sigma_{qQ}\right].
\end{equation*}
乘回颜色因子得到
\begin{equation*}
 \boxed{
 \overline{|\mathcal M|^2}
 =g_s^4\frac{N_c^2-1}{2N_c^2}
 \frac{(t-\Sigma_{qQ})^2+(u-\Sigma_{qQ})^2+2s\Sigma_{qQ}}{s^2},}
\end{equation*}
即正文\eqnrefs{eq:qqbar-to-QQbar-M2-dp10}。外腿质量壳同时给
\begin{equation*}
 s+t+u=2\Sigma_{qQ}.
\end{equation*}
最后沿 $m_q,m_Q\to0$ 的连续质量壳极限得到
\begin{equation*}
 \lim_{m_q,m_Q\to0}\overline{|\mathcal M|^2}
 =g_s^4\frac{N_c^2-1}{2N_c^2}\frac{t^2+u^2}{s^2}.
\end{equation*}
无质量 Born 核因此是有限质量树级结果的严格连续极限。
\end{derivation}



\subsection{$q\bar q\to gg$ 与 Ward 抵消}

QED 的 Compton 散射中，两种费米子线排序相加已经足以满足 Ward 恒等式。QCD 的 $q\bar q\to gg$ 还包含 $s$ 道三胶子图，因为两种夸克交换图的颜色算符分别按 $T^aT^b$ 与 $T^bT^a$ 排列；它们的差不是零，而是
\begin{equation*}
 T^aT^b-T^bT^a
 =[T^a,T^b]
 =\ii f^{abc}T^c.
\end{equation*}
因此仅保留两张夸克交换图时，纵向收缩留下一个系数为 $\ii f^{abc}T^c$ 的费米子流。三胶子顶角含同一个结构常数，其纵向收缩给出大小相同、符号相反的项。把三张树图写成
\begin{equation}
\boxed{
 \mathcal M_{ab}^{\mu\nu}
 =\mathcal M_{t,ab}^{\mu\nu}
 +\mathcal M_{u,ab}^{\mu\nu}
 +\mathcal M_{s,ab}^{\mu\nu}.}
\label{eq:qcd-qqbar-gg-amplitude-sum-dp68}
\end{equation}
对任一外胶子作纵向替换后，完整和满足
\begin{equation}
 \boxed{
 k_{1\mu}\mathcal M^{\mu\nu}_{ab}=0,
 \qquad
 k_{2\nu}\mathcal M^{\mu\nu}_{ab}=0.}
 \label{eq:qcd-qqbar-gg-ward-dp32}
\end{equation}
这条关系表明胶子自相互作用由非阿贝尔规范不变性固定，并在树级 $S$ 矩阵中保证非物理纵向偏振相消。

不同味夸克散射把一般 $SU(N_c)$ 的颜色平均压缩成 QCD 中熟悉的 $4/9$；交叉关系给出夸克--反夸克湮灭核；电磁产生中的末态颜色求和把 $N_c$ 直接写进包容 $R$ 比。$q\bar q\to gg$ 则进一步表明，非阿贝尔颜色算符的不对易性与三胶子自相互作用是同一规范结构的两个方面。圈图中，正是这些规范场和鬼场自由度共同改变强耦合的尺度运行并产生渐近自由。


相同味夸克散射比不同味过程多出一个不可回避的新结构：末态两个夸克完全相同，因此把两条外费米子线互换后得到的 $u$ 道图与原来的 $t$ 道图必须相干相加，而且 Fermi 统计给出相对负号。于是
\begin{equation}
\boxed{
\mathcal M(qq\to qq)=\mathcal M_t-\mathcal M_u,}
\label{eq:qcd-identical-qq-amplitude-dp95}
\end{equation}
其中 $\mathcal M_t$ 就是\eqnrefs{eq:qqprime-amplitude-dp10} 的同味版本，$\mathcal M_u$ 由 $p_3\leftrightarrow p_4$ 即 $t\leftrightarrow u$ 得到。对一般 $SU(N_c)$、无质量夸克，初态自旋和颜色平均后的结果为
\begin{equation}
\boxed{
\overline{|\mathcal M(qq\to qq)|^2}
=g_s^4\left\{
\frac{N_c^2-1}{2N_c^2}
\left[\frac{s^2+u^2}{t^2}+\frac{s^2+t^2}{u^2}\right]
-\frac{N_c^2-1}{N_c^3}\frac{s^2}{tu}
\right\}.}
\label{eq:qcd-identical-qq-m2-dp95}
\end{equation}
前两个平方项只是不同味结果及其交换像；第三项才是相同费米子交换的干涉。它同时依赖 Dirac 线的交换和颜色矩阵的非平凡收缩，不能通过把两个 Rutherford 型概率直接相加得到。物理 QCD 取 $N_c=3$ 后
\begin{equation*}
\overline{|\mathcal M|^2}
=g_s^4\left[
\frac49\frac{s^2+u^2}{t^2}
+\frac49\frac{s^2+t^2}{u^2}
-\frac{8}{27}\frac{s^2}{tu}
\right].
\end{equation*}

同一个四费米子解析振幅还把 $qq\to qq$ 与 $q\bar q\to q\bar q$ 连接起来。把一条外夸克从末态交叉到初态，相应四动量反号并把 $u$ 型自旋子换成 $v$ 型自旋子，Mandelstam 不变量随之置换。对相同味无质量 $q\bar q\to q\bar q$，树级同时存在 $s$ 道湮灭与 $t$ 道散射，结果为
\begin{equation}
\boxed{
\overline{|\mathcal M(q\bar q\to q\bar q)|^2}
=g_s^4\left\{
\frac{N_c^2-1}{2N_c^2}
\left[\frac{s^2+u^2}{t^2}+\frac{t^2+u^2}{s^2}\right]
-\frac{N_c^2-1}{N_c^3}\frac{u^2}{st}
\right\}.}
\label{eq:qcd-identical-qqbar-m2-dp95}
\end{equation}
这说明“交叉对称”不是把截面公式中的字母机械互换，而是先对解析振幅做外腿交叉，再重新进行物理通道的自旋、颜色平均；相同粒子统计和颜色干涉必须一起保留。

把\eqnrefs{eq:qcd-identical-qq-m2-dp95} 与 DP94 中的 $qg\to qg$、$gg\to gg$ 合在一起，QCD 的树级硬散射已经显出一个统一结构：运动学奇点由传播子通道决定，颜色系数由表示和生成元收缩决定，而完全相同外态还要求加入量子统计的交换干涉。后续的 PDF 因子化和喷注可观测量，就是把这些短程核嵌入强子初态并与软--共线长程结构分离。


前面的 $q\bar q\to gg$ 已经说明：一旦外态含胶子，非阿贝尔三胶子顶角不再是可有可无的修正，而是 Ward 恒等式所要求的树级组成部分。要把强子对撞机中最常见的 Born 硬核补齐，还需要再看两个含初态胶子的通道 $qg\to qg$ 与 $gg\to gg$。它们的意义不只是“再多算两个过程”：前者同时包含基本表示与伴随表示的颜色代数，后者则完全由胶子自相互作用组成，因而把 $C_F$ 与 $C_A$ 在可观测散射核中的不同角色直接显露出来。

先保留夸克质量 $m_q$，取
\begin{equation*}
 q(p_1,i)+g(k_1,a)\to q(p_2,j)+g(k_2,b),
 \qquad
 p_1^2=p_2^2=m_q^2c^2,
 \qquad k_1^2=k_2^2=0,
\end{equation*}
并定义
\begin{equation*}
 s=(p_1+k_1)^2,
 \qquad
 t=(p_1-p_2)^2,
 \qquad
 u=(p_1-k_2)^2,
 \qquad s+t+u=2m_q^2c^2.
\end{equation*}
树级振幅由 $s$、$u$ 道夸克交换以及 $t$ 道胶子交换三部分组成。把外胶子偏振矢量记为 $\varepsilon_\mu(k_1)$ 与 $\varepsilon_\nu^*(k_2)$，则可写成
\begin{equation}
\boxed{
\begin{aligned}
\mathcal M_{qg}(m_q)
=g_s^2\,\bar u(p_2)\Bigg[&
 T^bT^a\,\gamma^\nu
 \frac{\slashed p_1+\slashed k_1+m_qc}{s-m_q^2c^2}\gamma^\mu
 +T^aT^b\,\gamma^\mu
 \frac{\slashed p_1-\slashed k_2+m_qc}{u-m_q^2c^2}\gamma^\nu\\
&+\frac{\ii f^{abc}T^c}{t}\,
 \Gamma^{\rho\mu\nu}(p_1-p_2,k_1,-k_2)\gamma_\rho
\Bigg]u(p_1)\,
\varepsilon_\mu(k_1)\varepsilon_\nu^*(k_2).
\end{aligned}}
\label{eq:qcd-qg-amplitude-dp94}
\end{equation}
其中 $\Gamma^{\rho\mu\nu}$ 是前文已经由 $F^a_{\mu\nu}F^{a\mu\nu}$ 推出的三胶子顶角张量。第一条腿属于 $t$ 道内部胶子；因为三胶子规则把三条动量都定义为流入顶角，所以它携带的是 $p_1-p_2=k_2-k_1$，并满足 $(p_1-p_2)+k_1-k_2=0$。这条动量和是使用顶角公式前必须先做的检查。三项的相对系数由同一个协变导数与场强张量固定；若去掉第三项，纵向偏振替换不会相消。

有限 $m_q$ 的振幅由\eqnrefs{eq:qcd-qg-amplitude-dp94} 完整定义。强子对撞机的轻 parton Born 核进一步取严格连续极限 $m_q\to0$；此时 $s+t+u=0$。对初态夸克的 $N_c$ 个颜色和两个自旋、初态胶子的 $N_c^2-1$ 个颜色和两个物理偏振平均，并对末态求和，得到无质量部分子极限的闭式
\begin{equation}
\boxed{
\overline{|\mathcal M(qg\to qg)|^2}
=g_s^4(s^2+u^2)
\left[
 \frac1{t^2}-\frac{C_F}{C_A}\frac1{su}
\right],
\qquad
C_F=\frac{N_c^2-1}{2N_c},\quad C_A=N_c.}
\label{eq:qcd-qg-M2-dp94}
\end{equation}
这里第一项是前向 $t$ 道胶子交换产生的强增强，第二项来自夸克交换与非阿贝尔颜色排序之间的相干组合。对 QCD 的 $N_c=3$，$C_F/C_A=4/9$，恢复常用的
\begin{equation*}
 \overline{|\mathcal M|^2}
 =g_s^4(s^2+u^2)
 \left(\frac1{t^2}-\frac{4}{9su}\right).
\end{equation*}
由于物理区 $s>0$、$t<0$、$u<0$，第二项实际上为正；写成 $-1/(su)$ 的形式是为了保持交叉对称的 Mandelstam 结构。

纯胶子散射
\begin{equation*}
 g(k_1,a)+g(k_2,b)\to g(k_3,c)+g(k_4,d)
\end{equation*}
包含 $s,t,u$ 三个三胶子交换图以及一个四胶子接触图。用 Jacobi 恒等式把接触项分配到三个交换道后，完整树级振幅可按伴随表示颜色张量写成
\begin{equation}
\boxed{
 \mathcal M^{abcd}
 =g_s^2\left[
 f^{abe}f^{cde}\mathcal K_s
 +f^{ace}f^{bde}\mathcal K_t
 +f^{ade}f^{bce}\mathcal K_u
 \right],}
\label{eq:qcd-gg-color-amplitude-dp94}
\end{equation}
其中 $\mathcal K_s,\mathcal K_t,\mathcal K_u$ 包含相应交换图与四胶子接触项的洛伦兹结构。单独任何一项都依赖规范选择；只有它们的完整和在外腿上满足 Ward 恒等式。对一般 $SU(N_c)$ 完成颜色、偏振求和与初态平均后得到
\begin{equation}
\boxed{
\overline{|\mathcal M(gg\to gg)|^2}
=\frac{4g_s^4C_A^2}{N_c^2-1}
\left(
3-\frac{su}{t^2}-\frac{st}{u^2}-\frac{tu}{s^2}
\right).}
\label{eq:qcd-gg-M2-dp94}
\end{equation}
对 $N_c=3$，整体颜色系数成为 $9/2$。与 $qq'$ 或 $qg$ 相比，\eqnrefs{eq:qcd-gg-M2-dp94} 同时在 $t\to0$ 与 $u\to0$ 出现交换增强，反映了两个完全相同的无质量规范玻色子既可沿前向也可沿后向交换。式中的 $C_A^2$ 还直接表明：胶子自相互作用的强度由伴随表示 Casimir 控制，而不是由夸克的 $C_F$ 控制。

对于任意无质量 $2\to2$ Born 过程，统一相空间关系为
\begin{equation}
 \frac{\dd\sigma}{\dd t}
 =\frac{\hbar^2}{16\pi s^2}\overline{|\mathcal M|^2}.
 \label{eq:qcd-partonic-dsdt-master-dp94}
\end{equation}
因此 \eqnrefs{eq:qcd-qg-M2-dp94,eq:qcd-gg-M2-dp94} 已经足以直接生成角分布。\figref{fig:qcd-qg-gg-angular-dp94} 把 $N_c=3$ 时的无量纲 Born 核除以各自在 $90^\circ$ 的值；这样既去掉共同的 $g_s^4$ 与 $1/s$ 尺度，又保留由交换道和颜色结构决定的角形状。

\begin{figure}[htbp]
 \centering
 \includegraphics[width=.74\textwidth]{qcd_qg_gg_angular_dp94.pdf}
 \caption{QCD 中 $qg\to qg$ 与 $gg\to gg$ 的树级无量纲角分布核，均按 $90^\circ$ 归一化。$qg$ 通道主要显示 $t$ 道前向增强；$gg$ 通道还因相同外胶子而出现明显的后向交换增强。曲线由\eqnrefs{eq:qcd-qg-M2-dp94,eq:qcd-gg-M2-dp94} 逐点计算。}
 \label{fig:qcd-qg-gg-angular-dp94}
\end{figure}

\begin{derivation}[从\eqnrefs{eq:qcd-qg-amplitude-dp94}到\eqnrefs{eq:qcd-qg-M2-dp94}]
先取轻夸克极限 $m_q\to0$。这一极限下三张图的相对系数仍由规范不变作用量固定。单张图本身具有规范依赖，因此偏振求和先保留两种横向物理偏振；只有完整振幅相加并满足 Ward 恒等式以后，才可以使用相应的协变化简。

先固定质心系，使每个四动量和物理偏振都能逐分量检查。
令 $Q\equiv\sqrt{s}/2$，并取
\begin{align*}
 p_1^\mu&=Q(1,0,0,1),&
 k_1^\mu&=Q(1,0,0,-1),\\
 p_2^\mu&=Q(1,\sin\theta,0,\cos\theta),&
 k_2^\mu&=Q(1,-\sin\theta,0,-\cos\theta).
\end{align*}
于是
\begin{equation*}
 s=4Q^2,
 \qquad
 t=-2Q^2(1-\cos\theta),
 \qquad
 u=-2Q^2(1+\cos\theta),
 \qquad s+t+u=0.
\end{equation*}
记 $c_\theta\equiv\cos\theta$。入射胶子沿 $-z$ 方向，可选两条实的横向偏振
\begin{equation*}
 \varepsilon_{1x}^\mu=(0,1,0,0),
 \qquad
 \varepsilon_{1y}^\mu=(0,0,1,0),
\end{equation*}
而出射胶子沿 $(-\sin\theta,0,-\cos\theta)$，可选
\begin{equation*}
 \varepsilon_{2\parallel}^\mu=(0,-c_\theta,0,\sin\theta),
 \qquad
 \varepsilon_{2y}^\mu=(0,0,1,0).
\end{equation*}
逐项有 $k_i\cdot\varepsilon_{i\lambda}=0$ 与 $\varepsilon_{i\lambda}^2=-1$。因此这四种偏振组合恰好穷尽两条外胶子的物理态，不含时间样或纵向态。

三张图的 Dirac 核分别定义为
\begin{align*}
 \mathscr A_s^{\mu\nu}
 &=\gamma^\nu\frac{\slashed p_1+\slashed k_1}{s}\gamma^\mu,\\
 \mathscr A_u^{\mu\nu}
 &=\gamma^\mu\frac{\slashed p_1-\slashed k_2}{u}\gamma^\nu,\\
 \mathscr A_t^{\mu\nu}
 &=\frac1t\Gamma^{\rho\mu\nu}(p_1-p_2,k_1,-k_2)\gamma_\rho.
\end{align*}
这里第一条三胶子腿的动量必须是 $p_1-p_2=k_2-k_1$，因为前文三胶子顶角约定三条动量全部流入顶角：
\begin{equation*}
 (p_1-p_2)+k_1-k_2=0.
\end{equation*}
这也是正文\eqnrefs{eq:qcd-qg-amplitude-dp94} 中动量方向的直接检查。

接着计算六个独立的 Dirac--偏振收缩。
对 $x,y\in\{s,u,t\}$ 定义
\begin{equation*}
 K_{xy}(\theta)
 \equiv
 \sum_{r_1,r_2}
 \sum_{\lambda_1,\lambda_2}
 \Bigl[
 \bar u_{r_2}(p_2)\mathscr A_x^{\mu\nu}u_{r_1}(p_1)
 \varepsilon_{1\lambda_1,\mu}
 \varepsilon_{2\lambda_2,\nu}
 \Bigr]
 \Bigl[
 \bar u_{r_2}(p_2)\mathscr A_y^{\alpha\beta}u_{r_1}(p_1)
 \varepsilon_{1\lambda_1,\alpha}
 \varepsilon_{2\lambda_2,\beta}
 \Bigr]^* .
\end{equation*}
利用无质量自旋和
\begin{equation*}
 \sum_r u_r(p)\bar u_r(p)=\slashed p
\end{equation*}
后，每个 $K_{xy}$ 都只是 $4\times4$ Dirac 矩阵的有限迹。含三胶子顶角的项还可以先把 Lorentz 三指标降成一个普通 Dirac 矩阵。令 $q=p_1-p_2=k_2-k_1$，则
\begin{align*}
 &\varepsilon_1^\mu\varepsilon_2^\nu
 \Gamma^{\rho}{}_{\mu\nu}(q,k_1,-k_2)\\
 &\qquad={}
 -2(k_1\!\cdot\varepsilon_2)\varepsilon_1^\rho
 +(\varepsilon_1\!\cdot\varepsilon_2)(k_1+k_2)^\rho
 -2(k_2\!\cdot\varepsilon_1)\varepsilon_2^\rho.
\end{align*}
这一步直接来自三胶子顶角的三项，并只用了
$k_1\cdot\varepsilon_1=k_2\cdot\varepsilon_2=0$ 与 $q=k_2-k_1$。对上面选定的四种物理偏振，它进一步化为
\begin{align*}
 \mathscr A_t^{(x,y)}
 &=-\frac{2Q\sin\theta}{t}\slashed\varepsilon_y,
 &
 \mathscr A_t^{(y,\parallel)}
 &=-\frac{2Q\sin\theta}{t}\slashed\varepsilon_y,\\
 \mathscr A_t^{(y,y)}
 &=-\frac{\slashed k_1+\slashed k_2}{t},
 &
 \mathscr A_t^{(x,\parallel)}
 &=\frac{-2Q\sin\theta\slashed\varepsilon_{1x}
 +c_\theta(\slashed k_1+\slashed k_2)
 -2Q\sin\theta\slashed\varepsilon_{2\parallel}}{t}.
\end{align*}
因此 $t$ 道也已被化成普通的有限伽马矩阵链，不再剩下未计算的三胶子张量。作为第一项示范，$s$ 道平方变成
\begin{align*}
 K_{ss}
 ={}&\sum_{\lambda_1,\lambda_2}
 \Tr\!\left[
 \slashed p_2
 \slashed\varepsilon_2
 \frac{\slashed p_1+\slashed k_1}{s}
 \slashed\varepsilon_1
 \slashed p_1
 \slashed\varepsilon_1
 \frac{\slashed p_1+\slashed k_1}{s}
 \slashed\varepsilon_2
 \right].
\end{align*}
把 $p_1+k_1=(2Q,0,0,0)$ 与上面的四条物理偏振逐一代入，并反复使用
\begin{equation*}
 \{\gamma^\mu,\gamma^\nu\}=2\eta^{\mu\nu}I_4,
 \qquad
 \Tr(\slashed a\slashed b\slashed c\slashed d)
 =4\bigl(a\!\cdot b\,c\!\cdot d-a\!\cdot c\,b\!\cdot d+a\!\cdot d\,b\!\cdot c\bigr),
\end{equation*}
四种偏振组合相加给
\begin{equation*}
 K_{ss}=4(1+c_\theta).
\end{equation*}

再显式算一个含三胶子顶角的干涉项 $K_{st}$。定义 Dirac 共轭矩阵
\begin{equation*}
 \overline{\mathscr A}_t
 \equiv\gamma^0\mathscr A_t^\dagger\gamma^0,
\end{equation*}
则固定一组物理偏振 $(\lambda_1,\lambda_2)$ 时的自旋和为
\begin{equation*}
 K_{st}^{(\lambda_1,\lambda_2)}
 =\Tr\!\left[
 \slashed p_2\,
 \mathscr A_s^{(\lambda_1,\lambda_2)}
 \slashed p_1\,
 \overline{\mathscr A}_t^{(\lambda_1,\lambda_2)}
 \right].
\end{equation*}
把上一段已经化简出的四个 $\mathscr A_t$ 与
$\mathscr A_s=\slashed\varepsilon_2(\slashed p_1+\slashed k_1)\slashed\varepsilon_1/s$
逐一代入，四组偏振分别给
\begin{align*}
 K_{st}^{(x,\parallel)}
 &=\frac{4(1+c_\theta)}{1-c_\theta},&
 K_{st}^{(x,y)}
 &=2(1+c_\theta),\\
 K_{st}^{(y,\parallel)}
 &=2(1+c_\theta),&
 K_{st}^{(y,y)}
 &=\frac{4(1+c_\theta)}{1-c_\theta}.
\end{align*}
这里前、后两项中的 $1/(1-c_\theta)$ 来自 $t=-2Q^2(1-c_\theta)$；分子中的 $Q^2$ 则由两条外动量的 Dirac 迹提供，所以能量尺度消掉，只留下无量纲角函数。四项相加便得到
\begin{align*}
 K_{st}
 &=2(1+c_\theta)
 +2(1+c_\theta)
 +\frac{8(1+c_\theta)}{1-c_\theta}\\
 &=\frac{4(1+c_\theta)(3-c_\theta)}{1-c_\theta}.
\end{align*}
夸克传播子图与三胶子交换图的干涉因此归结为普通 Dirac 迹。物理干涉项还要乘颜色 Gram 元素 $G_{st}<0$，所以 $K_{st}$ 的符号本身不能判断最终干涉是增强还是相消。

剩余四个独立收缩使用同一套迹恒等式得到
\begin{align*}
 K_{uu}&=\frac{8(1+c_\theta^2)}{1+c_\theta},&
 K_{tt}&=\frac{8(1+c_\theta)(c_\theta^2-2c_\theta+5)}{(1-c_\theta)^2},\\
 K_{su}&=4(1-c_\theta),&
 K_{st}&=\frac{4(1+c_\theta)(3-c_\theta)}{1-c_\theta},&
 K_{ut}&=\frac{8(c_\theta^2-c_\theta+2)}{1-c_\theta}.
\end{align*}
由于 $K_{yx}=K_{xy}^*$，这六项已经覆盖三张图平方后的全部九个平方与干涉项。单个 $K_{xy}$ 依赖所选物理偏振基；与颜色结构相干组合并对完整物理振幅求和后，最终结果不依赖这一基底选择。

Lorentz--Dirac 收缩确定后，再独立计算颜色 Gram 矩阵。
把三张图的颜色矩阵记为
\begin{equation*}
 \mathcal C_s=T^bT^a,
 \qquad
 \mathcal C_u=T^aT^b,
 \qquad
 \mathcal C_t=\ii f^{abc}T^c=[T^a,T^b].
\end{equation*}
注意 $\mathcal C_t^\dagger=-\mathcal C_t$，因此干涉项的符号不能靠“颜色因子平方”猜测。定义
\begin{equation*}
 G_{xy}\equiv\sum_{a,b}\Tr_F(\mathcal C_x\mathcal C_y^\dagger).
\end{equation*}
由
\begin{equation*}
 \sum_aT^aT^a=C_FI,
 \qquad
 T^aT^bT^a=\left(C_F-\frac{C_A}{2}\right)T^b,
 \qquad
 \Tr(T^aT^b)=T_F\delta^{ab},
\end{equation*}
逐项得到
\begin{align*}
 G_{ss}=G_{uu}&=T_Fd_AC_F,\\
 G_{su}=G_{us}&=T_Fd_A\left(C_F-\frac{C_A}{2}\right),\\
 G_{tt}&=T_Fd_AC_A,\\
 G_{st}=G_{ts}&=-\frac12T_Fd_AC_A,\\
 G_{ut}=G_{tu}&=+\frac12T_Fd_AC_A,
\end{align*}
其中 $d_A=N_c^2-1$、$T_F=1/2$。例如
\begin{align*}
 G_{st}
 &=\sum_{a,b}\Tr\!\left[T^bT^a(T^bT^a-T^aT^b)\right]\\
 &=T_Fd_A\left(C_F-\frac{C_A}{2}\right)-T_Fd_AC_F
 =-\frac12T_Fd_AC_A,
\end{align*}
负的 $G_{st}$ 与正的 $G_{ut}$ 直接显示非阿贝尔交换图与两种夸克排序具有不同的颜色干涉符号。

最后先在振幅层相干求和，再实施初态平均。
未平均的完整平方振幅为
\begin{align*}
 \sum_{\rm spin,pol,color}|\mathcal M|^2
 =g_s^4\Bigl[&
 G_{ss}K_{ss}+G_{uu}K_{uu}+G_{tt}K_{tt}
 +2G_{su}K_{su}+2G_{st}K_{st}+2G_{ut}K_{ut}
 \Bigr].
\end{align*}
把上面六个 $K_{xy}$ 与六个颜色系数代入，并用
\begin{equation*}
 C_F=\frac{N_c^2-1}{2N_c},
 \qquad C_A=N_c,
 \qquad d_A=N_c^2-1,
\end{equation*}
整理成
\begin{equation*}
 \sum_{\rm spin,pol,color}|\mathcal M|^2
 =g_s^4\frac{d_A}{N_c}
 \frac{(c_\theta^2+2c_\theta+5)
 \left[N_c^2(c_\theta^2+2c_\theta+5)-(1-c_\theta)^2\right]}
 {(1-c_\theta)^2(1+c_\theta)}.
\end{equation*}
初态有 $2N_c$ 个夸克自旋--颜色态和 $2d_A$ 个胶子物理偏振--颜色态，所以平均因子是 $1/(4N_cd_A)$。再使用
\begin{align*}
 s^2+u^2&=4Q^4(c_\theta^2+2c_\theta+5),\\
 t^2&=4Q^4(1-c_\theta)^2,\\
 su&=-8Q^4(1+c_\theta),
\end{align*}
就得到
\begin{equation*}
 \overline{|\mathcal M(qg\to qg)|^2}
 =g_s^4(s^2+u^2)
 \left[
 \frac1{t^2}-\frac{C_F}{C_A}\frac1{su}
 \right],
\end{equation*}
即正文\eqnrefs{eq:qcd-qg-M2-dp94}。这里 $1/t^2$ 与 $1/(su)$ 是六类运动学收缩和六类颜色 Gram 元素相干组合后的结果；它们不能分别归属于某一张 Feynman 图的“概率”。
\end{derivation}

\begin{derivation}[从\eqnrefs{eq:qcd-gg-color-amplitude-dp94}到\eqnrefs{eq:qcd-gg-M2-dp94}]
这一过程若直接用 $-\eta_{\mu\nu}$ 展开四条外腿的偏振求和，会产生很长的 Lorentz 指标收缩，而且容易把“偏振代数”和“颜色代数”混在一起。这里改用物理螺旋度基底：先把每个非零螺旋度振幅算成三个颜色张量的线性组合，再单独完成一般 $SU(N_c)$ 的颜色求和。这样每一步都只处理一种数学结构。

先固定质心系运动学与外胶子的物理偏振。
令 $Q\equiv\sqrt{s}/2$，取度规 $\eta_{\mu\nu}=\operatorname{diag}(1,-1,-1,-1)$，并写
\begin{align*}
 k_1^\mu&=Q(1,0,0,1),&
 k_2^\mu&=Q(1,0,0,-1),\\
 k_3^\mu&=Q(1,\sin\theta,0,\cos\theta),&
 k_4^\mu&=Q(1,-\sin\theta,0,-\cos\theta).
\end{align*}
于是
\begin{equation*}
 t=-\frac{s}{2}(1-\cos\theta),
 \qquad
 u=-\frac{s}{2}(1+\cos\theta),
 \qquad s+t+u=0.
\end{equation*}
四个外胶子都只取横向物理偏振。以右螺旋度为例，可选
\begin{align*}
 \varepsilon_R^\mu(k_1)&=\frac1{\sqrt2}(0,1,\ii,0),&
 \varepsilon_R^\mu(k_2)&=\frac1{\sqrt2}(0,-1,\ii,0),\\
 \varepsilon_R^{*\mu}(k_3)&=\frac1{\sqrt2}(0,-\cos\theta,\ii,\sin\theta),&
 \varepsilon_R^{*\mu}(k_4)&=\frac1{\sqrt2}(0,\cos\theta,\ii,-\sin\theta).
\end{align*}
左螺旋度由上式作复共轭得到。逐一检查可得
\begin{equation*}
 k_i\cdot\varepsilon_{R,L}(k_i)=0,
 \qquad
 \varepsilon_{R,L}^*(k_i)\cdot\varepsilon_{R,L}(k_i)=-1,
\end{equation*}
所以这里从一开始就没有引入纵向或时间样偏振。

为缩短颜色指标，定义
\begin{equation*}
 \mathcal C_s\equiv f^{abe}f^{cde},
 \qquad
 \mathcal C_t\equiv f^{ace}f^{bde},
 \qquad
 \mathcal C_u\equiv f^{ade}f^{bce}.
\end{equation*}
Jacobi 恒等式给出
\begin{equation*}
 \mathcal C_s-\mathcal C_t+\mathcal C_u=0.
\end{equation*}

随后逐图计算独立的螺旋度振幅。
先算 $g_Rg_R\to g_Rg_R$。为避免与整体系数的无关相位纠缠，下表列出费曼规则给出的 $\ii\mathcal M$。三行交换图分别含两个三胶子顶角和一个 $1/s$、$1/t$、$1/u$ 传播子，最后一行为四胶子接触顶角。把上面的四个偏振矢量逐分量代入三胶子顶角
\begin{equation*}
 \Gamma^{\mu\nu\rho}(p,q,r)
 =\eta^{\mu\nu}(p-q)^\rho
 +\eta^{\nu\rho}(q-r)^\mu
 +\eta^{\rho\mu}(r-p)^\nu,
 \qquad p+q+r=0,
\end{equation*}
即可得到
\begin{center}
\begin{tabular}{c|c}
图 & $\ii\mathcal M(g_Rg_R\to g_Rg_R)/(\ii g_s^2)$\\ \hline
$s$ 道 & $-\mathcal C_s\cos\theta$\\[2pt]
$t$ 道 & $\displaystyle
 \mathcal C_t\frac{19+7\cos\theta-11\cos^2\theta+\cos^3\theta}
 {4(1-\cos\theta)}$\\[8pt]
$u$ 道 & $\displaystyle
 \mathcal C_u\frac{19-7\cos\theta-11\cos^2\theta-\cos^3\theta}
 {4(1+\cos\theta)}$\\[8pt]
四胶子接触项 & $\displaystyle
 -\mathcal C_s\cos\theta
 -\frac{\mathcal C_t}{4}(3+2\cos\theta-\cos^2\theta)
 -\frac{\mathcal C_u}{4}(3-2\cos\theta-\cos^2\theta)$
\end{tabular}
\end{center}
这里每一行只是有限个四维分量的乘加。以 $s$ 道为例，传播子分母为 $(k_1+k_2)^2=s$；两个三胶子顶角分别与
$\varepsilon_R(k_1)\varepsilon_R(k_2)$ 和
$\varepsilon_R^*(k_3)\varepsilon_R^*(k_4)$ 收缩。利用
$k_i\cdot\varepsilon_i=0$ 后，分子中的 $Q^2=s/4$ 恰好抵消传播子的 $1/s$，留下表中的 $-\mathcal C_s\cos\theta$。$t$、$u$ 道保留的 $1/(1\mp\cos\theta)$ 正是后来前向、后向交换极点的来源；接触项没有传播子，因此只给角度多项式。

把四行相加，并用 $\mathcal C_s=\mathcal C_t-\mathcal C_u$，分子多项式发生大规模相消：
\begin{align*}
 \frac{\ii\mathcal M_{RR\to RR}}{\ii g_s^2}
 &=4\left(
 \frac{\mathcal C_t}{1-\cos\theta}
 +\frac{\mathcal C_u}{1+\cos\theta}
 \right)\\
 &=-2\left(
 \mathcal C_t\frac{s}{t}
 +\mathcal C_u\frac{s}{u}
 \right).
\end{align*}
这一步显示了四胶子接触项的作用：它本身没有交换极点，却与三个交换图的多项式部分相消，使完整振幅只剩满足 Ward 恒等式与交叉结构的组合。

其余非零螺旋度振幅由同一分量计算或交叉变换得到。取三个独立通道可写成
\begin{align*}
 \mathcal M_{RR\to RR}
 &=-2g_s^2s\left(\frac{\mathcal C_t}{t}+\frac{\mathcal C_u}{u}\right),\\
 \mathcal M_{LR\to LR}
 &=2g_s^2u\left(\frac{\mathcal C_t}{t}+\frac{\mathcal C_s}{s}\right),\\
 \mathcal M_{LR\to RL}
 &=-2g_s^2t\left(\frac{\mathcal C_s}{s}-\frac{\mathcal C_u}{u}\right).
\end{align*}
整体相位不影响平方振幅。宇称变换给出另外三个等模振幅；其余螺旋度组合在树级为零。因此螺旋度求和只需把这三个独立结果平方后乘以 $2$；因子 $2$ 来自宇称相关的另外三个等模通道。

螺旋度核确定后，计算一般 $SU(N_c)$ 的颜色 Gram 矩阵。
记伴随表示维数
\begin{equation*}
 d_A=N_c^2-1.
\end{equation*}
首先，三个颜色张量的平方范数相同。例如
\begin{align*}
 \sum_{a,b,c,d}\mathcal C_s^2
 &=\sum_{e,e'}
 \left(\sum_{a,b}f^{abe}f^{ab e'}\right)
 \left(\sum_{c,d}f^{cde}f^{cd e'}\right)\\
 &=C_A^2\sum_{e,e'}\delta^{ee'}\delta^{ee'}
 =C_A^2d_A.
\end{align*}
令伴随表示生成元为 $(F^a)_{bc}=-\ii f^{abc}$。由
\begin{equation*}
 \Tr(F^aF^b)=C_A\delta^{ab},
 \qquad
 \sum_aF^aF^bF^a=\frac{C_A}{2}F^b,
\end{equation*}
可得一个代表性的交叉颜色积
\begin{align*}
 \sum_{a,b,c,d}\mathcal C_s\mathcal C_t
 &=\sum_{e,e'}\Tr(F^eF^{e'}F^eF^{e'})\\
 &=\frac{C_A}{2}\sum_{e'}\Tr(F^{e'}F^{e'})
 =\frac12C_A^2d_A.
\end{align*}
再结合 Jacobi 关系 $\mathcal C_s-\mathcal C_t+\mathcal C_u=0$ 及外颜色指标的置换对称性，得到完整 Gram 矩阵
\begin{equation*}
 G_{xy}\equiv\sum_{a,b,c,d}\mathcal C_x\mathcal C_y
 =C_A^2d_A
 \begin{pmatrix}
 1&\tfrac12&-\tfrac12\\
 \tfrac12&1&\tfrac12\\
 -\tfrac12&\tfrac12&1
 \end{pmatrix}_{xy},
 \qquad x,y\in\{s,t,u\}.
\end{equation*}
负的 $G_{su}$ 并不表示负概率；它只记录两个颜色张量在颜色向量空间中的夹角，最终概率由完整振幅的 Gram 二次型给出。

最后完成颜色求和、螺旋度求和与初态平均。
利用上面的 Gram 矩阵，第一个独立螺旋度通道给出
\begin{align*}
 \sum_{\rm color}|\mathcal M_{RR\to RR}|^2
 &=4g_s^4C_A^2d_A\,s^2
 \left(
 \frac1{t^2}+\frac1{u^2}+\frac1{tu}
 \right).
\end{align*}
同理
\begin{align*}
 \sum_{\rm color}|\mathcal M_{LR\to LR}|^2
 &=4g_s^4C_A^2d_A\,u^2
 \left(
 \frac1{t^2}+\frac1{s^2}+\frac1{ts}
 \right),\\
 \sum_{\rm color}|\mathcal M_{LR\to RL}|^2
 &=4g_s^4C_A^2d_A\,t^2
 \left(
 \frac1{s^2}+\frac1{u^2}+\frac1{su}
 \right).
\end{align*}
宇称伙伴再给相同贡献，因此全螺旋度、全颜色求和为
\begin{equation*}
 \sum_{\rm hel,color}|\mathcal M|^2
 =8g_s^4C_A^2d_A\,\mathcal S(s,t,u),
\end{equation*}
其中
\begin{align*}
 \mathcal S(s,t,u)={}&
 s^2\left(\frac1{t^2}+\frac1{u^2}+\frac1{tu}\right)
 +t^2\left(\frac1{s^2}+\frac1{u^2}+\frac1{su}\right)\\
 &+u^2\left(\frac1{s^2}+\frac1{t^2}+\frac1{st}\right).
\end{align*}
利用 $s+t+u=0$，例如
\begin{equation*}
 -\frac{2su}{t^2}
 =\frac{s^2+u^2}{t^2}-1,
\end{equation*}
以及循环置换，可把它化成
\begin{equation*}
 \mathcal S(s,t,u)
 =2\left(
 3-\frac{su}{t^2}-\frac{st}{u^2}-\frac{tu}{s^2}
 \right).
\end{equation*}
每个初态胶子有 $d_A$ 个颜色和两个物理螺旋度，所以初态简并度为 $4d_A^2$。最终
\begin{align*}
 \overline{|\mathcal M|^2}
 &=\frac{1}{4d_A^2}
 \sum_{\rm hel,color}|\mathcal M|^2\\
 &=\frac{4g_s^4C_A^2}{d_A}
 \left(
 3-\frac{su}{t^2}-\frac{st}{u^2}-\frac{tu}{s^2}
 \right),
\end{align*}
再代入 $d_A=N_c^2-1$，即得到正文\eqnrefs{eq:qcd-gg-M2-dp94}。
\end{derivation}


\begin{derivation}[从\eqnrefs{eq:qcd-identical-qq-amplitude-dp95}到\eqnrefs{eq:qcd-identical-qq-m2-dp95}]
把两张图写为
\begin{align*}
\mathcal M_t&=\frac{g_s^2}{t}
[\bar u_3\gamma^\mu T^a u_1]
[\bar u_4\gamma_\mu T^a u_2],\\
\mathcal M_u&=\frac{g_s^2}{u}
[\bar u_4\gamma^\mu T^a u_1]
[\bar u_3\gamma_\mu T^a u_2].
\end{align*}
相同末态费米子交换给 $\mathcal M=\mathcal M_t-\mathcal M_u$，所以
\begin{equation*}
\overline{|\mathcal M|^2}
=\overline{|\mathcal M_t|^2}
+\overline{|\mathcal M_u|^2}
-2\Re\overline{\mathcal M_t\mathcal M_u^*}.
\end{equation*}
前两项直接由不同味结果得到
\begin{equation*}
\overline{|\mathcal M_t|^2}
=g_s^4\frac{N_c^2-1}{2N_c^2}\frac{s^2+u^2}{t^2},
\qquad
\overline{|\mathcal M_u|^2}
=g_s^4\frac{N_c^2-1}{2N_c^2}\frac{s^2+t^2}{u^2}.
\end{equation*}
真正的新信息在交叉项。颜色平均可写成一个单迹：
\begin{align*}
\mathcal C_{tu}
&=\frac1{N_c^2}\sum_{a,b}\Tr(T^aT^bT^aT^b).
\end{align*}
利用
\begin{equation*}
T^aT^bT^a=\left(C_F-\frac{C_A}{2}\right)T^b
=-\frac{1}{2N_c}T^b,
\end{equation*}
以及 $\sum_b\Tr(T^bT^b)=T_F(N_c^2-1)$，得到
\begin{equation*}
\mathcal C_{tu}
=-\frac{N_c^2-1}{4N_c^3}.
\end{equation*}
自旋交叉项不能拆成两个独立迹，而是形成一个八 $\gamma$ 链。先按完备关系把四条外腿的自旋求和，注意复共轭振幅会反转第二条费米子链的矩阵次序，于是
\begin{equation*}
\mathcal S_{tu}
=\frac14\Tr\!\left(
\slashed p_3\gamma^\mu\slashed p_1\gamma^\nu
\slashed p_4\gamma_\mu\slashed p_2\gamma_\nu
\right).
\end{equation*}
这里不能把迹“按视觉上相邻的四个矩阵”拆开。先把中间五个 $\gamma$ 矩阵作为整体收缩。四维 Clifford 代数给出
\begin{equation*}
\gamma^\mu\gamma^\alpha\gamma^\nu\gamma^\beta\gamma_\mu
=-2\gamma^\beta\gamma^\nu\gamma^\alpha,
\end{equation*}
因此
\begin{align*}
\mathcal S_{tu}
&=-\frac12\Tr\!\left(
\slashed p_3\slashed p_4\gamma^\nu
\slashed p_1\slashed p_2\gamma_\nu
\right).
\end{align*}
第二次收缩只剩四个普通四矢量。利用
\begin{equation*}
\gamma^\nu\slashed a\slashed b\gamma_\nu=4a\!\cdot b,
\qquad
\Tr(\slashed a\slashed b)=4a\!\cdot b,
\end{equation*}
逐步得到
\begin{align*}
\mathcal S_{tu}
&=-\frac12\,4(p_1\!\cdot p_2)\Tr(\slashed p_3\slashed p_4)\\
&=-8(p_1\!\cdot p_2)(p_3\!\cdot p_4).
\end{align*}
对无质量 $2\to2$ 运动学，$2p_1\!\cdot p_2=s$ 且 $2p_3\!\cdot p_4=s$，故
\begin{equation*}
\boxed{\mathcal S_{tu}=-2s^2.}
\end{equation*}
这一步说明干涉项只留下 $s^2$ 并不是由 $s+t+u=0$ 猜出的，而是八 $\gamma$ 链在 Clifford 代数下真正收缩后的结果。
于是
\begin{equation*}
-2\Re\overline{\mathcal M_t\mathcal M_u^*}
=-2\frac{g_s^4}{tu}\,\mathcal C_{tu}\mathcal S_{tu}
=-g_s^4\frac{N_c^2-1}{N_c^3}\frac{s^2}{tu},
\end{equation*}
与两个平方项相加即得到\eqnrefs{eq:qcd-identical-qq-m2-dp95}。注意颜色迹本身为负，而 Fermi 交换又给振幅相对负号；最终干涉符号必须在这两层结构全部保留后才能确定。
\end{derivation}

\begin{derivation}[从\eqnrefs{eq:qcd-qqbar-gg-amplitude-sum-dp68}到\eqnrefs{eq:qcd-qqbar-gg-ward-dp32}]
把两条出射胶子动量记为 $k_1,k_2$，颜色指标为 $a,b$。将三张图共同的 $g_s^2$ 因子提出，并定义约化张量
\begin{equation*}
 \widetilde{\mathcal M}_x^{\mu\nu}
 \equiv g_s^{-2}\mathcal M_x^{\mu\nu},
 \qquad x=t,u,s.
\end{equation*}
对两张夸克交换图，费米子线结构严格写成
\begin{align*}
 \widetilde{\mathcal M}_t^{\mu\nu}
 &=\bar v(p_2)\gamma^\nu T^b
 S(p_1-k_1)\gamma^\mu T^a u(p_1),\\
 \widetilde{\mathcal M}_u^{\mu\nu}
 &=\bar v(p_2)\gamma^\mu T^a
 S(p_1-k_2)\gamma^\nu T^b u(p_1),
\end{align*}
其中 $S(q)$ 表示 Dirac 算符 $S^{-1}(q)=\slashed q-mc$ 的 Feynman 逆核。代数运算先在内部费米子分母非零的开运动学区域进行，所得恒等式再按统一的 $+\ii0$ 处方延拓。

把第一条胶子偏振替换为 $k_1^\mu$。由
\begin{equation*}
 \slashed k_1
 =S^{-1}(p_1)-S^{-1}(p_1-k_1)
\end{equation*}
以及在壳方程 $S^{-1}(p_1)u(p_1)=0$，得到
\begin{align*}
 S(p_1-k_1)\slashed k_1u(p_1)
 &=S(p_1-k_1)
 [S^{-1}(p_1)-S^{-1}(p_1-k_1)]u(p_1)\\
 &=-u(p_1).
\end{align*}
因此 $k_{1\mu}\widetilde{\mathcal M}_t^{\mu\nu}$ 留下一个没有内部传播子的费米子流，颜色顺序为 $T^bT^a$。

对另一种排序，利用四动量守恒 $p_1+p_2=k_1+k_2$，把 $k_1$ 写成靠近反费米子端的两个逆传播子之差，并令逆传播子向左作用到 $\bar v(p_2)$。由 $\bar v(p_2)(\slashed p_2+mc)=0$，得到同一个旋量流，颜色顺序为 $T^aT^b$，符号与上一项相反。两张约化夸克图的纵向收缩之和因而含有
\begin{equation*}
 T^aT^b-T^bT^a=[T^a,T^b]=\ii f^{abc}T^c
\end{equation*}
乘同一个费米子流。阿贝尔理论的生成元是普通数，所以这一对易子严格为零；QCD 的非零对易子要求三胶子图参与同一 Ward 恒等式。

$s$ 道图含夸克颜色流
\begin{equation*}
 J_c^\rho=\bar v(p_2)\gamma^\rho T^c u(p_1)
\end{equation*}
以及由 Yang--Mills 三次作用量得到的三胶子张量 $V_{\rho}{}^{\mu\nu}$。按“全部动量流入”的统一约定收缩 $k_{1\mu}V^{\rho\mu\nu}$，再使用 $k_1^2=k_2^2=0$、$k_2\cdot\varepsilon_2=0$ 与 $q=k_1+k_2$，得到
\begin{equation*}
 k_{1\mu}\widetilde{\mathcal M}_s^{\mu\nu}
 =-k_{1\mu}
 \left(\widetilde{\mathcal M}_t^{\mu\nu}
 +\widetilde{\mathcal M}_u^{\mu\nu}\right).
\end{equation*}
恢复共同因子 $g_s^2$ 后，三图满足
\begin{equation*}
 k_{1\mu}
 \left(\mathcal M_t^{\mu\nu}
 +\mathcal M_u^{\mu\nu}
 +\mathcal M_s^{\mu\nu}\right)=0.
\end{equation*}
交换 $1\leftrightarrow2$ 同样得到第二条 Ward 恒等式。三图抵消依次使用 Dirac 逆传播子的望远镜恒等式、颜色生成元的对易关系和由同一 Yang--Mills 作用量固定的三胶子顶角。
\end{derivation}

